% ===== 原文 PDF pp.38--40（印刷 pp.37--39）=====
% 译者 tr-03；图 5 归属 tr-03
% [原文 p.37]
\begin{center}
  \includegraphics[width=0.62\linewidth]{fig5.png}
\end{center}
\figtitle{5}{库仑分支上特殊坐标 $(a,a_D)$ 的不同选取，以及绕奇点的对偶单值性。}

% [Tutor Note: 图 5 中四个圆盘是 $\cM_{\mathrm{CB}}$（右下角标注 $u$-plane / $\cM_{\mathrm{CB}}$）
% 上互相重叠的开片，各开片内标出该处的特殊坐标二重态：$(a,a^D)$、$(a',a'^D)$、$(a'',a''^D)$、
% $(a''',a'''^D)$；红色菱形是一条绕中心奇点（$\times$）的闭合路径，它穿过四个开片、每穿过一次
% 就做一次 $\SL(2,\Z)$ 对偶变换，回到出发点时给出单值性 $\cM_\gamma$（见下文 (3.26)）。
% 注意图内标注用上标 $a^D$，而图注用下标 $a_D$——这是原文自身的记号不一致，不是图注排错。]

我们可以平等地对待这两个坐标。也就是说，既然坐标 $a$ 表示 $\cN=2$ 矢量多重态中的标量分量，
我们就可以把另一个坐标 $a^D$ 视为 $a$ 的磁对偶，于是 $a^D$ 属于一个包含对偶规范场 $A^D$ 的
$\cN=2$ 矢量多重态。以 $a^D$ 表述的对偶描述可以由 $a$ 与预势 $\cF(a)$ 通过 (3.16) 确定，因此
两种描述携带相同的信息，而这些信息都被 $\cF$ 所捕获。从一个开片 $\cU$（二重态
$(a(u),a^D(u))$）走到另一个开片 $\cU'$（另一个二重态 $(a'(u),a'^D(u))$）会涉及一个
$\SL(2,\Z)$ 对偶变换。的确，二重态 $(a,a^D)$ 可以看作奇点之外的模空间
$\cM_{\mathrm{CB}}$ 上某个 $\SL(2,\Z)$ 丛的全纯截面，并绕 $\cM_{\mathrm{CB}}$ 上的奇点经历
非平凡的 $\SL(2,\Z)$ 单值性（更多细节见 3.2 节）。这基本上就是电磁对偶如何编码在 CB 上的
要点。图 5 给出了应记住的示意图。

类似的分析可以推广到更高秩 $r$ 的理论，此时 $\cM_{\mathrm{CB}}$ 是 $r$ 复维的空间。对偶群
变为 $\Sp(2r,\Z)$，而每个特殊坐标 $\mathbf a$ 与 $\mathbf a_D$ 现在都可以看作 $r$ 维矢量，
其变换为
\begin{equation}
  \begin{pmatrix} \mathbf a^{\mathbf D}\\[2pt] \mathbf a\end{pmatrix}\to
  \begin{pmatrix} A & B\\[2pt] C & D\end{pmatrix}
  \begin{pmatrix} \mathbf a^{\mathbf D}\\[2pt] \mathbf a\end{pmatrix},
  \tag{3.24}
\end{equation}
其中 $A,B,C,D$ 是 $r\times r$ 矩阵，它们一起参数化 $\Sp(2r,\Z)$ 群。

% [原文 p.38]
\SUB{3.2}{库仑分支上的预势与特殊几何}

正如上一节所解释的，CB 有效场论由 CB 各开片上的局域描述拼成，这些开片通过电磁对偶变换
粘合在一起，而低能物理完全编码在全纯预势 $\cF(a)$ 之中。本节我们简要介绍求解 $\cF(a)$ 的
一般策略。确定了预势，就可以找出式 (3.23) 中描述有效场论库仑分支的特殊坐标。

首先，对预势有几条一般性约束：
\begin{itemize}
  \item 它必须是 $a$ 的局域全纯函数；
  \item 它必须尊重 UV 理论中存在的对称性，例如 $\mathrm{U}(1)_r$ R 对称性；
  \item 最后，在大 VEV 极限 $a\to\infty$ 下，CB 上的物理是弱耦合的，因此我们可以信任直接
        使用 UV 拉格朗日量所做的计算。这意味着，在这些区域，预势必须与能从 UV 理论得到的
        结果相容。
\end{itemize}

在这些一般性约束下，预势的一般表达式由 [148] 给出
\begin{equation}
  \resizebox{0.82\textwidth}{!}{$\displaystyle
  \cF(a) = \frac{1}{2}\tau_{\mathrm{UV}}a^2
  + \frac{ib}{8\pi}a^2\ln\left(\frac{a^2}{\Lambda^2}\right)
  + \sum_{k=1}^{\infty}F_k\left(\frac{\Lambda}{a}\right)^{bk}a^2\,,$}
  \tag{3.25}
\end{equation}
这是弱耦合区域 $|a|>\Lambda$ 中有效的展开，其中 $\Lambda$ 是动力学产生的标度。第二项来自
CB 弱耦合极限下的微扰一圈修正，即式 (2.40) 的一圈跑动，而这一贡献中的系数 $b$ 已在式
(2.37) 与 (2.41) 中定义。第三项来自非微扰修正，即可能的瞬子修正。系数 $F_k$ 可以用不同方式
确定：

\begin{enumerate}
  \item 一种方式是直接计算以 $e^{2\pi n i\tau}$ 加权的瞬子效应。这一技术已对 $\SU(N)$
        （以及 $\mathrm{U}(N)$）规范群在 [76,77] 中发展出来，随后推广到其他经典规范群，如
        $\SO(N)$ 与 $\Sp(N)$ [159,160]。\footnote{在我们的约定下 $\Sp(1)\simeq\SU(2)$。}
  \item 或者，也可以利用 $F(a)$ 全纯这一事实，通过它在 $\cM_{\mathrm{CB}}$ 上各奇点附近的
        行为来确定它，这正是 [20,21] 中所展开的故事的主旨。
\end{enumerate}

% [Tutor Note: (3.25) 三项依次是：树阶项 $\frac12\tau_{\mathrm{UV}}a^2$、一圈对数修正
% $\frac{ib}{8\pi}a^2\ln(a^2/\Lambda^2)$、以及瞬子求和。注意求和的幂次是 $bk$ 而不是 $k$：
% $b$ 就是紧接下文所说的 (2.37)、(2.41) 中的一圈系数，故第 $k$ 阶瞬子修正携带
% $(\Lambda/a)^{bk}$ 的标度幂。展开参数的“小”由 $|a|>\Lambda$ 保证。]

% [Tutor Note: 原文在此处把预势写成 $F(a)$（斜体 $F$）而非别处的花体 $\cF(a)$，第二项“利用
% $F(a)$ 全纯”的那一条同样如此；已用 900 dpi 裁图复核，确认是原文的排版不一致而非文字层误差。
% 本译本按原样复刻，读者应把它与 $\cF$ 视为同一个预势。]

% [原文 p.39]
让我们把第二种方法的更多细节讲清楚。$\cM_{\mathrm{CB}}$ 中的奇点\footnote{按 [63] 的说法，
这里的奇点是指 Kähler 度规 $g_{ij}$ 出现奇点，而不是 $\cM_{\mathrm{CB}}$ 所携带的复结构出现
奇点。后者具有不同的物理含义，相关研究见 [161,162]。}是特殊坐标 $(a,a_D)$ 不再单值的地方，
因此库仑分支有效场论在那里失效。\footnote{除尖点型之外，奇点处的有效规范耦合取有限值
（例如见 [163]）。}一个简单的例子是 $\SU(2)$ SCFT 中位于原点 $u=0$ 处的情形 [21]，那里
$a=\sqrt{u/2}$。这正是 $a$ 不能在 $\cM_{\mathrm{CB}}$ 上定义整体坐标的原因。物理上，CB 上
出现这些奇点，是因为某些带电粒子在这些奇点处变得无质量；若把它们积分掉，就会导致有效场论
中出现奇点 [20,21]。

这些奇点在特殊几何中有一种用单值性给出的解释。当我们在 $u$ 平面上沿一条环路 $\gamma$ 绕过
奇点 $s$ 时，特殊坐标 $(a^D,a)$ 会获得某个单值性 $\cM_\gamma$，使得
\begin{equation}
  \begin{pmatrix} a^D\\[2pt] a\end{pmatrix}\to
  \cM_\gamma\begin{pmatrix} a^D\\[2pt] a\end{pmatrix}.
  \tag{3.26}
\end{equation}
因为 $a^D$ 是 $a$ 的函数，若没有奇点，就不会有非平凡的单值性。单值性刻画了 CB 上的奇点。
此外，这些单值性受到
\begin{equation}
  \prod_{\{\gamma\}}\cM_\gamma = \cM_\infty\,,
  \tag{3.27}
\end{equation}
的约束，意思是：与 CB 上各奇点相关联的所有单值性矩阵之积，必须等于由 $u$ 平面中环绕所有
奇点（在无穷远处）的路径给出的单值性矩阵。而这正是瞬子修正可以忽略的区域，因此
$\cM_\infty$ 可以由 UV 理论的一圈计算确定。于是式 (3.27) 约束了 CB 上奇点的可能结构，并且
事实证明，它对式 (3.25) 中的系数 $F_k$ 施加了极其有力的约束 [20,21]，使得只需少量额外假设，
就能一并确定所有 $F_k$，而无需做任何显式的瞬子计算。关于该论证的更多细节，我们建议读者
参阅 [20,21]（或例如综述 [84,164]）。

% [Tutor Note: (3.27) 左端的乘积遍历所有把奇点在 $u$ 平面中围起来的环路 $\{\gamma\}$，右端
% $\cM_\infty$ 对应一条在“无穷远”处把所有奇点围在内部的大环路。由于每个 $\cM_\gamma$ 都由
% 一条带基点的环路决定，乘积的次序依赖环路的组合次序；原文只用 $\prod$ 略去了这一细节。
% 该式也解释了为何 $F_k$ 能被唯一确定：奇点附近的局部单值性由曲线退化决定，而整体乘积由
% 弱耦合的一圈计算固定，两者相匹配即给出约束。]

% [Tutor Note: 本块末（p.40 页末）为完整句，句号收尾；PDF p.41 首行是“Example 3.1:”，由 tr-04
% 接手。§3.2 的正文自此继续。]

% [边界] p.40 页末为完整句；下一位 tr-04 自 p.41 首行（例 3.1）继续，无需补写本块末句。
